\documentclass[10pt,a4paper]{amsart}
\usepackage[margin=19mm]{geometry}
\usepackage{amsmath,amssymb,mathtools}
\usepackage[scr=rsfs,cal=euler]{mathalfa}
\usepackage{xfrac}
\usepackage[unicode,hidelinks,bookmarks=false]{hyperref}
\newcommand{\ie}{i.e.\ }
\newcommand{\cf}{cf.\ }
\newcommand{\resp}{respectively\ }
\newcommand{\Subsection}{\subsection}
\providecommand{\nobreakdash}{\nobreak\hbox{-}}
\setlength{\emergencystretch}{2em}
\multlinegap0pt
\usepackage{bm}
\usepackage{bbm}
\usepackage{dsfont}
\usepackage{xspace}
\usepackage{cancel}
\usepackage{mathtools}
\usepackage{amsthm}
\makeatletter
\@ifundefined{thm}{\newtheorem{thm}{Thm}}{}
\@ifundefined{prop}{\newtheorem{prop}[thm]{Proposition}}{}
\makeatother
\DeclareMathOperator{\Ext}{Ext}
\DeclareMathOperator{\Sym}{Sym}
\DeclareMathOperator{\ct}{ct}
\newcommand{\A}{\mathcal{A}}
\newcommand{\M}{\mathcal{M}}
\newcommand{\Tor}{\mathsf{Tor}}
\renewcommand{\O}{\mathcal{O}}
\title[Tautological and non-tautological cycles on the moduli space]{Tautological and non-tautological cycles on the moduli space of abelian varieties}
\author{Samir Canning, Dragos Oprea, Rahul Pandharipande}
\date{Source: arXiv:2408.08718v3 (CC BY 4.0)}
\begin{document}
\maketitle

\noindent\emph{Source and licence.} Tautological and non-tautological cycles on the moduli space of abelian varieties. Version arXiv:2408.08718v3 (CC BY 4.0), distributed under the Creative Commons Attribution 4.0 International licence, as recorded in the arXiv licence metadata for this exact version and shown on the abstract page https://arxiv.org/abs/2408.08718v3. Full rights notice: licenses/arxiv-2408.08718-CC-BY-4.0.txt.\par\smallskip

\noindent\textit{Editorial scope.} Counted unit: the Prop whose source label is \texttt{strictns} in the article (source0.tex lines 1258--1260, source label \texttt{strictns}), delivered together with the article's own complete original proof of it (source0.tex lines 1261--1287, one \texttt{proof} environment of 27 lines). No mathematics is written, repaired or re-derived: statement and proof are the article's own text line for line. Required context retained uncounted: none. Every definition, display and label of those lines is retained unchanged. Omitted from this package and not redistributed: 1--1257, 1288--3449. Nothing inside a retained excerpt is omitted or altered: every retained body is delivered exactly as the article's lines are written, so no omission receipt of any kind accompanies this package. Citations: the retained excerpts carry no citation command at all, so no bibliography entry of the article is transcribed and no reference is redistributed. Reference adaptation: none was needed and none was made; every reference inside the delivered unit resolves inside it, so no unresolved reference or citation is produced. Printed slips kept literal: no printed word, symbol, hypothesis or number of the counted statement or of its proof was altered. Layout and class: the article's own class is \texttt{amsart} with packages including tikz-cd, enumitem, amssymb, amsmath, amscd, mathrsfs, url, cite, fullpage, hyperref, verbatim, comment, fancyvrb, fvextra; the wrapper is the lane's pinned \texttt{amsart} editorial wrapper at 10pt on A4, which restates the theorem-like environments the retained text needs and adds the article's own macro definitions that the retained text needs (\texttt{\textbackslash A}, \texttt{\textbackslash Ext}, \texttt{\textbackslash M}, \texttt{\textbackslash O}, \texttt{\textbackslash Sym}, \texttt{\textbackslash Tor}, \texttt{\textbackslash ct}), with extra wrapper packages (bm, bbm, dsfont, xspace, cancel, mathtools). The delivered numbering is the unit's own. Assets: no retained span contains a figure environment, a graphics-inclusion command, a PDF-page inclusion, a table or a TikZ picture; no image file is copied into this package, the package declares no asset dependency and no image file is redistributed with it. Licence: version arXiv:2408.08718v3 of this article is distributed under the Creative Commons Attribution 4.0 International licence, as recorded in the arXiv licence metadata for this exact version and shown on the abstract page https://arxiv.org/abs/2408.08718v3. A full rights notice is delivered at licenses/arxiv-2408.08718-CC-BY-4.0.txt. Characters: every retained excerpt is pure ASCII, so the delivered engine, XeLaTeX with its default Latin Modern text font, reports no missing character.
\begin{prop} \label{strictns}
    The stack theoretic image of $\epsilon^{\circ}_{\mathsf{I}}$ is nonsingular. In particular, $\Tor^{-1}(\A_1\times \A_{g-1})$ is nonsingular away from the intersection of its components. 
\end{prop}

\begin{proof}
    %By definition, we have a factorization
    %\[
    %\M^{\circ \ct}_{\mathsf{I}}\rightarrow \Tor^{-1}(\A_1\times \A_{g-1})\rightarrow \M^{\ct}_{g}
    %\]$\epsilon^{\circ}_{\mathsf{I}}$ is a bijection onto its image. Because $\M^{\circ \ct}_{\mathsf{I}}$ is nonsingular, Zariski's main theorem implies that $\epsilon_{\mathsf{I}}^{\circ}$ is an isomorphism onto $\epsilon_{\mathsf{I}}^{\circ}(\M^{\circ \ct}_{\mathsf{I}})$ if $\epsilon_{\mathsf{I}}^{\circ}(\M^{\circ \ct}_{\mathsf{I}})$ is nonsingular. 

    The tangent space of $\Tor^{-1}(\A_1\times \A_{g-1})$ at a point $(C,(E, B))$ is the fiber product of the tangent space $\Ext^1(\Omega_C,\O_C)$ to $\M_g^{\ct}$ at $C$ with the tangent space $\Sym^2 H^0(\Omega_E)^{\vee}\oplus \Sym^2 H^0(\Omega_B)^{\vee}$ to $\A_1\times \A_{g-1}$ at $(E,B)$ over the tangent space $$\Sym^2 H^0(\Omega_E)^{\vee}\oplus (H^0(\Omega_E)^{\vee}\otimes H^0(\Omega_B)^{\vee})  \oplus \Sym^2 H^0(\Omega_B)^{\vee}$$  at $J(C)\cong E\times B$ of $\A_g$. 
    
    Let $k$ denote the number of leaves of the extremal tree $\mathsf{I}$. 
    Assume first $k=1$, for simplicity. A general point of $\epsilon_{\mathsf{I}}^{\circ}(\M^{\circ \ct}_{\mathsf{I}})$ is the Jacobian of a curve $C=E\cup D$, where $E$ is nonsingular of genus $1$ and $D$ is nonsingular of genus $g-1$. There is a tangent vector $v\in \Ext^1(\Omega_C,\O_C)$ corresponding to the smoothing of the node $p$ of $C$, hence $v\in T_pE \otimes T_pD$. %By \cite[Proposition 2.8]{FriedmanSmith}, 
    Under the differential of the Torelli map
    \[
    \Ext^1(\Omega_C,\O_C)\rightarrow \Sym^2 H^0(\omega_C)^{\vee}\cong \Sym^2 H^0(\Omega_E)^{\vee}\oplus (H^0(\Omega_E)^{\vee}\otimes H^0(\Omega_{D})^{\vee})  \oplus \Sym^2 H^0(\Omega_D)^{\vee}\,,
    \]
    $v$ maps to a nonzero vector in $(H^0(\Omega_E)^{\vee}\otimes H^0(\Omega_{D})^{\vee})$. Hence, $v$ does not lie in the tangent space to $\Tor^{-1}(\A_1\times \A_{g-1})$ at $(C, (E,J(D)))$, which thus has codimension at least $1$ in $\Ext^1(\Omega_C,\O_C)$. Because $\M^{\circ \ct}_{\mathsf{I}}$ is of codimension $1$ in $\M_g^{\ct}$, we see that $\epsilon_{\mathsf{I}}^{\circ}(\M^{\circ \ct}_{\mathsf{I}})$ is nonsingular at $(C,(E,J(D)))$. 
    
    Next, we suppose $C=E\cup D$, where $D=\cup_{i=0}^{n} D_i$ is a compact type curve of genus $g-1$ glued to $E$ at a exactly one point on $D_0$, where $D_0$ is nonsingular of genus $0<h<g-1$. Again, we consider the tangent vector $v$ corresponding to a family of curves smoothing the node $E\cap D_0$. Under the Torelli map, this family maps to $\A_{h+1}\times \A_{g-h-1}\subset \A_g$. %As in the proof of Proposition \ref{zeroint}, we identify the fiber product of $\A_{h-1}\times \A_{g-h+1}$ with $\A_1\times \A_{g-1}$ over $\A_g$ with $\A_1\times \A_{h-1}\times \A_{g-1-h}$. 
    Therefore, we can view the codomain of the differential of the Torelli map as 
    \[
    \Sym^2 H^0(\Omega_E)^{\vee}\oplus (H^0(\Omega_E)^{\vee}\otimes H^0(\Omega_{D_0})^{\vee}) \oplus \Sym^2 H^0(\Omega_{D_0})^{\vee}\oplus \Sym^2 \left(\bigoplus_{i=1}^{n} H^0(\Omega_{D_i})^{\vee}\right)\,,
    \]
    where the first three summands correspond to the $\A_{h+1}$ factor and the latter summands correspond to the $\A_{g-h-1}$ factor.
    As above, the tangent vector $v$ has nonzero image in the $(H^0(\Omega_E)^{\vee}\otimes H^0(\Omega_{D_0})^{\vee})$ summand. %by \cite[Proposition 2.8]{FriedmanSmith}. 
    Hence, the vector $v$ does not lie in the tangent space to $\Tor^{-1}(\A_1\times \A_{g-1})$ at $(C, (E,J(D)))$, and the conclusion follows as in the previous paragraph. 

    The cases when $k>1$ are proved analogously by analyzing the image of the tangent vectors corresponding to smoothings of the nodes represented by edges in $\mathsf{I}$.
     \end{proof}

\end{document}
